\section{Summary}

This document approaches the dynamics of disadvantageous traits from a biophysical perspective, treating populations as dynamical systems whose evolution is shaped by the interplay between sexual selection, mutation, and predation \cite{nowakEvolutionaryDynamicsExploring2006}. Within this framework, sexual selection is understood as a frequency-dependent process, in which the reproductive value of a trait depends on the traits and choices of other individuals in the population. The aim of this work was to
Develop a stochastic model of phenotype coexistence in a population under sexual selection and predation, based on the Moran process under weak selection, and embed a fitness landscape into the model. The model is instantiated for male guppies (\textit{Poecilia reticulata}), a system chosen because of the extensive empirical data available on sexual selection and predation in this species, which allows the theoretical hypotheses to be tested rather than merely illustrated. 


[Methods]: To this end, we [metodología breve: muestreo, experimento, análisis] muestreo: For these implementations, we used Guppies \textit{Poecilia reticulata} for the extended studies about them, female mate choice strongly favours colourful males \cite{houdeMATECHOICEBASED1987}, yet this same trait increases predation risk \cite{magnhagen1991, godinPredatorPreferenceBrightly2003}. 
 over [periodo o condiciones]. [Results]: Our results show that [hallazgo principal 1] , and that [hallazgo principal 2]. [Conclusion/Implication]: These findings suggest that [interpretación o implicación], providing insight into evolutionary dynamics for sexual selection. 


\begin{center}
    \textit{Keywords: Evolution, Evolutionary Game Theory (EGT), Sexual Selection, Stable Equilibrium States, Noise and Chaos, Moran Process, Poecilia Reticulata}
\end{center}

